\originalchapter{费曼图与散射振幅}\label{ch:feynman}
\originalsection{Toy theory 的定义}
取三个实标量场 $A,B,C$, 自由项为各自 $\tfrac12(\partial\phi)^2-\tfrac12m_\phi^2\phi^2$, 相互作用 $\mathcal L_{\rm int}=-gABC$. $g$ 在四维具有质量量纲. 规定顶点为 $-\ii g$, 内部线为 $\ii/(q^2-m^2+\ii0)$, 每个独立回路积 $\int\dd^4q/(2\pi)^4$, 所有图之和等于 $\ii\mathcal M$.

一个 $A\to B+C$ 顶点给 $\mathcal M=-g$. 树图没有自由回路动量, 顶点的能动量守恒把每条内部线动量确定. 内线称虚粒子, 一般不满足壳上条件; 不要用 $E^2=\mathbf p^2+m^2$ 把传播子分母错误设成零.

\originalsection{交换图与干涉}
$A+A\to B+B$ 由 C 交换, 有 $t$、$u$ 两图. 给
\begin{equation}
 \mathcal M=-g^2\left(\frac1{t-m_C^2}+\frac1{u-m_C^2}\right).
\end{equation}
末态的两种安排不可区分, 所以先把振幅相加再平方. $|M_t+M_u|^2$ 中有干涉项, 并不是两条过程概率的普通相加. 同粒子 $1/2!$ 是相空间计数, 不应同时删一张交换图和再除两次.

\begin{center}
\begin{tikzpicture}[x=0.8cm,y=0.8cm]
\coordinate(a)at(1,1);\coordinate(b)at(1,-1);\coordinate(c)at(3,1);\coordinate(d)at(3,-1);
\draw(0,2)node[left]{$A$}--(a)--(4,2)node[right]{$B$};\draw(0,-2)node[left]{$A$}--(b)--(4,-2)node[right]{$B$};\draw[dashed](a)--(b)node[midway,right]{$C$};
\node at(2,-2.5){$t$ 通道示意};
\end{tikzpicture}
\end{center}
图中线段只表示连接关系, 不表示真实粒子轨迹. 在强相互作用或量子电动力学中, 线型还标记自旋、颜色和粒子种类, 并附带矩阵顶点.

\originalsection{回路与正规化}
$A+A\to A+A$ 最低阶需要四个 ABC 顶点, 内部交替 B、C, 为一回路盒图. 四个带标记外腿在环上有三种不等价循环次序, 每种有 B、C 两种交替安排. 对质量均为零的 B、C, 两安排的积分相同, 但原图仍有两份.

若外腿全部按流入回路记为 $k_1,k_2,k_3,k_4$, 总和为零, 一个循环次序给
\begin{equation}
 \ii\mathcal M_{1234}=g^4\int\frac{\dd^4q}{(2\pi)^4}
 \frac1{(q^2+\ii0)[(q+k_1)^2+\ii0][(q+k_1+k_2)^2+\ii0][(q-k_4)^2+\ii0]}.
\end{equation}
全振幅为对三种次序与两种内部标记求和. $\ii0$ 指定时间排序的极点边界条件, 不是可随意忽略的实数小参数. 某些回路在紫外或红外发散, 需要正规化后与相应物理定义组合. ``积分形式已写出''与``积分值已有限求得''必须区分.

\originalsection{力程与传播子}
静态交换的空间势与 $1/(\mathbf q^2+m^2)$ 的三维 Fourier 变换相关. 沿极轴积分角变量后,
\begin{equation}
 \int\frac{\dd^3q}{(2\pi)^3}\frac{\e^{\ii\mathbf q\cdot\mathbf r}}{q^2+m^2}
 =\frac{\e^{-mr}}{4\pi r}.
\end{equation}
因此有质量媒介产生 Yukawa 型指数衰减, 力程约 $1/m=\hbar/(mc)$; 无质量光子产生长程 Coulomb 势. 这里指静态低能近似, 不能把 $1/m$ 当作虚粒子必须活过某个确定时间后消失的实际轨迹长度.
